\setcounter{chapter}{17}
\chapter{Containerization}\label{ch:18-containerization}

\chapterrule

The Notes API now has pinned dependencies, externalized configuration, a runtime contract, a clean directory structure, structured logging, and a build process with automated tests. But it still runs as a raw Python process, directly on whatever machine it happens to be installed on. The environment~--- the Python version, the operating system libraries, the directory layout~--- is assumed but not guaranteed.

This assumption is the root cause of environment drift (\hyperref[ch:07-environment-drift]{Chapter 7}), hidden dependencies (\hyperref[ch:08-hidden-dependencies]{Chapter 8}), and the ``it works on my machine'' problem (\hyperref[ch:09-configuration-coupling]{Chapter 9}). Even with careful management, the gap between a developer's laptop running macOS 14 with Python 3.11.8 and a production server running Debian 12 with Python 3.11.2 can cause subtle failures.

\textbf{Containerization} eliminates this gap by packaging the application together with its complete execution environment into a single, self-contained unit called a \textbf{container image}. The image includes:

\begin{itemize}
\tightlist
\item
  The operating system libraries
\item
  The exact Python version
\item
  All installed Python packages at their exact versions
\item
  The application code
\item
  The startup command
\end{itemize}

\noindent{}When this image runs~--- on a developer's laptop, on a CI server, on a production server in Singapore~--- it runs identically. The execution environment is identical because it travels with the code. The ``it works on my machine'' problem disappears because ``my machine'' and ``the server'' are now running exactly the same image. (One gap remains: CPU architecture. An image built on an Apple Silicon laptop is \texttt{arm64}, while most servers are \texttt{x86-64}; \texttt{docker\ build\ -}\allowbreak{}\texttt{-}\allowbreak{}\texttt{platform\ linux/}\allowbreak{}\texttt{amd64} builds for the server's architecture.)

\textbf{Docker} is the dominant tool for building and running container images. This chapter explains what containers actually are at the operating system level, builds a production-quality Docker image for the Notes API, and establishes the \texttt{docker\ build} + \texttt{docker\ run} workflow.

\begin{objectives}

By the end of this chapter, you will understand:

\begin{itemize}
\tightlist
\item
  What a Linux container is and how it differs from a virtual machine
\item
  What a Dockerfile is and how Docker uses it to build an image
\item
  How to write a production-quality Dockerfile for the Notes API
\item
  How Docker layers work and why they matter for build speed
\item
  How to run the containerized application with environment variables
\item
  How \texttt{docker\ compose} simplifies multi-container development (application + database)
\end{itemize}

\end{objectives}

\setcounter{section}{0}
\section{What a Container Is}\label{18-containerization:181-what-a-container-is}

Before writing any Docker files, it is important to understand what a container actually is. Many explanations describe containers with analogies (``like a lightweight VM''), but the accurate explanation is more precise and more useful.

\subsection*{The three Linux kernel mechanisms}\phantomsection\label{18-containerization:the-three-linux-kernel-mechanisms}

A Linux container is an ordinary operating system process with three specific kernel mechanisms applied:

\textbf{1. Namespaces~--- isolation}

A Linux \textbf{namespace} restricts what a process can see. There are different types:

\begin{itemize}
\tightlist
\item
  \textbf{PID namespace}: the container process sees its own isolated process ID tree. Inside the container, the first process has PID 1. From outside, it has a different PID.
\item
  \textbf{Network namespace}: the container has its own network interfaces, IP addresses, and routing table. Traffic is isolated from the host network and from other containers.
\item
  \textbf{Mount namespace}: the container sees its own filesystem. The root (\texttt{/}) inside the container is not the same as the root on the host.
\item
  \textbf{UTS namespace}: the container has its own hostname, different from the host's hostname.
\end{itemize}

\noindent{}Namespaces restrict visibility. A process inside a namespace cannot see processes outside it. (The host, in the parent namespace, can still see every container's processes~--- which is how \texttt{ps} on the server shows them.)

\textbf{2. cgroups~--- resource limits}

\textbf{Control groups (cgroups)} limit how much of a resource a process can consume:

\begin{itemize}
\tightlist
\item
  CPU: this container can use at most 0.5 CPU cores
\item
  Memory: this container can use at most 512 MB RAM
\item
  Disk I/O: this container can read/write at most 100 MB/s
\end{itemize}

\noindent{}Without cgroups, one container could consume all available memory and starve other containers. With cgroups, each container's consumption is bounded.

\textbf{3. A layered filesystem (overlayfs)}

The container's filesystem is built from \textbf{layers}, stacked on top of each other. Each layer is read-only. When a process writes a file, the write goes to a writable layer on top of the read-only layers. The read-only layers are shared between all containers running the same image~--- only the writable layer is unique per container.

\subsection*{Containers vs virtual machines}\phantomsection\label{18-containerization:containers-vs-virtual-machines}

A \textbf{virtual machine (VM)} virtualizes complete hardware: a \textbf{hypervisor}, helped by the CPU's virtualization features, gives each VM its own virtual CPUs, RAM, and disk. A full operating system (kernel + userspace) runs inside the VM. The overhead is noticeable: each VM runs its own kernel and system services, typically costing hundreds of megabytes of RAM before any application starts.

A \textbf{container} shares the host's kernel. No hardware emulation. No separate OS kernel. The container's ``OS'' is just userspace files (libraries, binaries)~--- not a running kernel. The overhead is minimal: a container starts in milliseconds and uses only the RAM that the application itself needs.

\begin{diagrambox}[short]{344.8pt}
\begin{Verbatim}[fontsize=\fontsize{9.0}{8.55}\selectfont]
VIRTUAL MACHINE                         CONTAINER
┌─────────────────────────────┐         ┌─────────────────────────────┐
│  Application                │         │  Application                │
│  Python + dependencies      │         │  Python + dependencies      │
│  OS userspace (glibc, etc.) │         │  OS userspace (glibc, etc.) │
│  Guest OS kernel            │         └─────────────────────────────┘
└─────────────────────────────┘         ┌─────────────────────────────┐
┌─────────────────────────────┐         │  HOST KERNEL (shared)       │
│  Hypervisor                 │         └─────────────────────────────┘
│  Host kernel and hardware   │         (no hypervisor, no guest kernel)
└─────────────────────────────┘

Startup time: seconds to a minute       Startup time: 100ms–1 second
Memory overhead: hundreds of MB         Memory overhead: a few MB
\end{Verbatim}
\end{diagrambox}

\noindent{}Containers are faster, lighter, and more portable than VMs. They are not more secure (sharing the kernel means a kernel vulnerability can escape the isolation), but for most application deployment purposes, containers provide adequate isolation with dramatically lower overhead.

\setcounter{section}{1}
\section{Docker: The Container Toolchain}\label{18-containerization:182-docker-the-container-toolchain}

\textbf{Docker} is the toolchain that makes containers practical for developers. It provides:

\begin{itemize}
\tightlist
\item
  \textbf{Docker daemon}: the background service that manages containers and images
\item
  \textbf{Docker CLI}: the \texttt{docker} command you use in your terminal
\item
  \textbf{Dockerfile}: a text file that describes how to build an image
\item
  \textbf{Docker Hub}: a public registry where images are stored and shared
\item
  \textbf{Docker Compose}: a tool for running multiple containers together
\end{itemize}

\subsection*{Installing Docker}\phantomsection\label{18-containerization:installing-docker}

On macOS and Windows, install \textbf{Docker Desktop} from \href{https://docker.com}{docker.com}. Docker Desktop includes the daemon, CLI, and Docker Compose.

On Linux, install Docker Engine:

\begin{codebox}[short]

\begin{Shaded}
\begin{Highlighting}[]
\CommentTok{\# Ubuntu/Debian}
\FunctionTok{sudo}\NormalTok{ apt{-}get update}
\FunctionTok{sudo}\NormalTok{ apt{-}get install docker.io docker{-}compose{-}v2}
\FunctionTok{sudo}\NormalTok{ usermod }\AttributeTok{{-}aG}\NormalTok{ docker }\VariableTok{$USER}  \CommentTok{\# Allow running docker without sudo}
\CommentTok{\# Log out and back in for group membership to take effect}
\end{Highlighting}
\end{Shaded}

\end{codebox}

\noindent{}Verify the installation:

\begin{codebox}[short]

\begin{Shaded}
\begin{Highlighting}[]
\ExtensionTok{docker} \AttributeTok{{-}{-}version}
\CommentTok{\# Docker version 25.0.3, build 4debf41}

\ExtensionTok{docker}\NormalTok{ run hello{-}world}
\CommentTok{\# Hello from Docker! ...}
\end{Highlighting}
\end{Shaded}

\end{codebox}

\setcounter{section}{2}
\section{The Dockerfile}\label{18-containerization:183-the-dockerfile}

A \textbf{Dockerfile} is a text file containing a sequence of instructions that Docker executes, one by one, to build an image. Each instruction that changes files (\texttt{RUN}, \texttt{COPY}, \texttt{ADD}) adds a layer to the image; the others (\texttt{WORKDIR}, \texttt{ENV}, \texttt{CMD}, \ldots) only record settings.

\subsection*{A minimal Dockerfile (what NOT to use in production)}\phantomsection\label{18-containerization:a-minimal-dockerfile-what-not-to-use-in-production}

\begin{codebox}[short]

\begin{Shaded}
\begin{Highlighting}[]
\CommentTok{\# DO NOT USE — this is a bad Dockerfile for illustration purposes}
\KeywordTok{FROM}\NormalTok{ python:3.11}
\KeywordTok{WORKDIR}\NormalTok{ /app}
\KeywordTok{COPY}\NormalTok{ . .}
\KeywordTok{RUN} \ExtensionTok{pip}\NormalTok{ install }\AttributeTok{{-}r}\NormalTok{ requirements.txt}
\KeywordTok{CMD}\NormalTok{ [}\StringTok{"python"}\NormalTok{, }\StringTok{"run.py"}\NormalTok{]}
\end{Highlighting}
\end{Shaded}

\end{codebox}

\noindent{}This works but has serious problems:

\begin{itemize}
\tightlist
\item
  The \texttt{python:3.11} base image is enormous (\textasciitilde1 GB)~--- it includes compilers and build tools that the application does not need at runtime
\item
  Copying \texttt{.} includes \texttt{.env}, \texttt{venv/}, \texttt{.git/}, and other files that should not be in the image
\item
  No non-root user~--- the application runs as root inside the container, which is a security risk
\item
  No separation between dependency installation and code copying~--- every code change rebuilds all layers including pip install
\end{itemize}

\subsection*{The production-quality Dockerfile}\phantomsection\label{18-containerization:the-production-quality-dockerfile}

\begin{codeboxcap}{\textcolor{accent}{Listing 18.1}\enspace Dockerfile (production-quality)}

\begin{Shaded}
\begin{Highlighting}[]
\CommentTok{\# =============================================================================}
\CommentTok{\# Multi{-}stage build for the Notes API}
\CommentTok{\#}
\CommentTok{\# A multi{-}stage build uses multiple FROM instructions. The first stage (builder)}
\CommentTok{\# installs the Python dependencies. The second stage (final) copies only what}
\CommentTok{\# is needed at runtime. Anything the builder needed only to *install* things}
\CommentTok{\# (compilers, headers, pip\textquotesingle{}s temporary files) never reaches the final image.}
\CommentTok{\# =============================================================================}

\CommentTok{\# ─────────────────────────────────────────────────────────────────────────────}
\CommentTok{\# Stage 1: builder}
\CommentTok{\# Install the Python dependencies into a separate directory.}
\CommentTok{\# ─────────────────────────────────────────────────────────────────────────────}
\KeywordTok{FROM}\NormalTok{ python:3.11{-}slim }\KeywordTok{AS}\NormalTok{ builder}

\CommentTok{\# No system packages are needed here: every dependency in requirements.txt}
\CommentTok{\# ships as a prebuilt wheel (psycopg2{-}binary bundles its own copy of libpq).}
\CommentTok{\# If a dependency ever has to be compiled from source (plain psycopg2, for}
\CommentTok{\# example), its build tools are installed in THIS stage:}
\CommentTok{\#   RUN apt{-}get update \&\& apt{-}get install {-}y {-}{-}no{-}install{-}recommends \textbackslash{}}
\CommentTok{\#       gcc libpq{-}dev \&\& rm {-}rf /var/lib/apt/lists/*}
\CommentTok{\# and the final stage then needs only the runtime library (libpq5).}

\CommentTok{\# Set the working directory for subsequent instructions in this stage.}
\CommentTok{\# Nothing in /build is copied to the final image, so it does not appear there.}
\KeywordTok{WORKDIR}\NormalTok{ /build}

\CommentTok{\# Copy ONLY the requirements file first.}
\CommentTok{\# Docker caches each layer. If requirements.txt has not changed,}
\CommentTok{\# Docker reuses the cached pip install layer from the previous build.}
\CommentTok{\# This means code changes (which are copied later) do not trigger}
\CommentTok{\# a full pip reinstall — only dependency changes do.}
\KeywordTok{COPY}\NormalTok{ requirements.txt .}

\CommentTok{\# Install Python dependencies into a known location.}
\CommentTok{\# {-}{-}no{-}cache{-}dir: do not cache downloaded packages (saves space)}
\CommentTok{\# {-}{-}prefix=/install: install into /install instead of system site{-}packages}
\CommentTok{\#   This makes it easy to copy just the packages to the final image.}
\KeywordTok{RUN} \ExtensionTok{pip}\NormalTok{ install }\AttributeTok{{-}{-}no{-}cache{-}dir} \AttributeTok{{-}{-}prefix}\OperatorTok{=}\NormalTok{/install }\AttributeTok{{-}r}\NormalTok{ requirements.txt}

\CommentTok{\# ─────────────────────────────────────────────────────────────────────────────}
\CommentTok{\# Stage 2: final}
\CommentTok{\# The actual container image that will run in production.}
\CommentTok{\# It contains the packages from stage 1, but none of stage 1\textquotesingle{}s leftovers.}
\CommentTok{\# ─────────────────────────────────────────────────────────────────────────────}
\KeywordTok{FROM}\NormalTok{ python:3.11{-}slim }\KeywordTok{AS}\NormalTok{ final}

\CommentTok{\# Create a non{-}root user for running the application.}
\CommentTok{\# Running as root inside a container is a security risk:}
\CommentTok{\# if the application is compromised, the attacker has root access}
\CommentTok{\# inside the container (and potentially on the host in some configurations).}
\CommentTok{\#}
\CommentTok{\# useradd options:}
\CommentTok{\#   {-}{-}create{-}home: create a home directory for the user}
\CommentTok{\#   {-}{-}shell /bin/bash: set the login shell}
\CommentTok{\#   {-}{-}uid 1001: set a specific UID (avoids conflicts with host UIDs)}
\CommentTok{\#   notes: the username}
\KeywordTok{RUN} \ExtensionTok{useradd} \AttributeTok{{-}{-}create{-}home} \AttributeTok{{-}{-}shell}\NormalTok{ /bin/bash }\AttributeTok{{-}{-}uid}\NormalTok{ 1001 notes}

\CommentTok{\# Copy the installed Python packages from the builder stage.}
\CommentTok{\# COPY {-}{-}from=builder copies files from the named stage.}
\CommentTok{\# /install contains all the packages installed by pip in stage 1.}
\CommentTok{\# /usr/local is where Python looks for packages by default.}
\KeywordTok{COPY} \OperatorTok{{-}{-}from=builder}\NormalTok{ /install /usr/local}

\CommentTok{\# Set the working directory for the application.}
\CommentTok{\# This is where the application code will live inside the container.}
\KeywordTok{WORKDIR}\NormalTok{ /app}

\CommentTok{\# WORKDIR creates /app owned by root. Hand the directory itself to \textquotesingle{}notes\textquotesingle{}}
\CommentTok{\# so the application can create files in it (a SQLite database, say).}
\CommentTok{\# This runs before the code is copied, so it stays cached.}
\KeywordTok{RUN} \FunctionTok{chown}\NormalTok{ notes:notes /app}

\CommentTok{\# Copy the application code, owned by the non{-}root user.}
\CommentTok{\# Note: this is after copying requirements/installing packages.}
\CommentTok{\# Changing application code only invalidates these layers,}
\CommentTok{\# not the pip install layer above — builds are fast.}
\CommentTok{\#}
\CommentTok{\# {-}{-}chown sets the owner as the files are copied. (A separate}
\CommentTok{\# "RUN chown {-}R notes:notes /app" afterwards would re{-}run on every code}
\CommentTok{\# change and store a second copy of every file in a new layer.)}
\CommentTok{\#}
\CommentTok{\# We use .dockerignore (below) to exclude files that should not be in the image.}
\KeywordTok{COPY} \OperatorTok{{-}{-}chown=notes:notes}\NormalTok{ app/ ./app/}
\KeywordTok{COPY} \OperatorTok{{-}{-}chown=notes:notes}\NormalTok{ run.py .}

\CommentTok{\# Switch to the non{-}root user.}
\CommentTok{\# All subsequent instructions and the running container use this user.}
\KeywordTok{USER}\NormalTok{ notes}

\CommentTok{\# Document which port the application listens on.}
\CommentTok{\# EXPOSE does not actually publish the port — it is documentation.}
\CommentTok{\# The port is published when running the container with {-}p or in docker{-}compose.}
\KeywordTok{EXPOSE}\NormalTok{ 5000}

\CommentTok{\# Set default environment variables.}
\CommentTok{\# These are the defaults; they can be overridden at runtime with {-}e flags.}
\CommentTok{\# Do NOT set SECRET\_KEY or DATABASE\_URL here — those are secrets and must}
\CommentTok{\# be provided at runtime, not baked into the image.}
\CommentTok{\# (HOST and PORT are not set: they are read only by run.py, and in the}
\CommentTok{\# container Gunicorn\textquotesingle{}s {-}{-}bind, below, decides the address and port.)}
\KeywordTok{ENV}\NormalTok{ FLASK\_DEBUG=false }\OperatorTok{\textbackslash{}}
\NormalTok{    LOG\_FORMAT=json }\OperatorTok{\textbackslash{}}
\NormalTok{    LOG\_LEVEL=INFO }\OperatorTok{\textbackslash{}}
\NormalTok{    PYTHONUNBUFFERED=1 }\OperatorTok{\textbackslash{}}
\NormalTok{    PYTHONDONTWRITEBYTECODE=1}

\CommentTok{\# PYTHONUNBUFFERED=1: Python outputs are sent straight to the terminal}
\CommentTok{\#   without buffering. Without this, log output might not appear until}
\CommentTok{\#   a buffer fills — which means last log entries before a crash might}
\CommentTok{\#   never appear.}
\CommentTok{\#}
\CommentTok{\# PYTHONDONTWRITEBYTECODE=1: Python does not write .pyc files inside the}
\CommentTok{\#   container. These are unnecessary in a container (the image is immutable)}
\CommentTok{\#   and waste space and inode counts.}

\CommentTok{\# Health check: Docker periodically runs this command to determine}
\CommentTok{\# if the container is healthy. If it fails repeatedly, Docker marks}
\CommentTok{\# the container as "unhealthy". Plain Docker only reports that status}
\CommentTok{\# (in docker ps); orchestrators such as Docker Swarm or Kubernetes act}
\CommentTok{\# on it by replacing the container.}
\CommentTok{\#}
\CommentTok{\# The slim image has no curl, so the check uses Python itself:}
\CommentTok{\# urlopen() raises an exception, and python exits non{-}zero, if the}
\CommentTok{\# request fails, times out, or returns an HTTP error status.}
\CommentTok{\# interval=30s: check every 30 seconds}
\CommentTok{\# timeout=5s: the command must complete within 5 seconds}
\CommentTok{\# retries=3: after 3 consecutive failures, mark as unhealthy}
\CommentTok{\# start{-}period=10s: give the container 10 seconds to start before checking}
\KeywordTok{HEALTHCHECK} \OperatorTok{{-}{-}interval=30s} \OperatorTok{{-}{-}timeout=5s} \OperatorTok{{-}{-}retries=3} \OperatorTok{{-}{-}start{-}period=10s}\NormalTok{ \textbackslash{}}
    \KeywordTok{CMD} \ExtensionTok{python} \AttributeTok{{-}c} \StringTok{"import urllib.request; urllib.request.urlopen(\textquotesingle{}http://localhost:5000/health\textquotesingle{}, timeout=3)"} \KeywordTok{||} \BuiltInTok{exit}\NormalTok{ 1}

\CommentTok{\# The command that runs when the container starts.}
\CommentTok{\#}
\CommentTok{\# We use Gunicorn instead of Flask\textquotesingle{}s built{-}in development server.}
\CommentTok{\# Gunicorn is a production WSGI server (Chapter 24) that:}
\CommentTok{\# {-} Runs several worker processes to serve requests in parallel}
\CommentTok{\#   (Flask\textquotesingle{}s development server is not built for production traffic)}
\CommentTok{\# {-} Restarts workers that crash or hang}
\CommentTok{\# {-} Handles graceful shutdown correctly}
\CommentTok{\#}
\CommentTok{\# Gunicorn imports the app package and calls create\_app() itself, so run.py}
\CommentTok{\# is not used here. The startup checks still run, because create\_app() runs}
\CommentTok{\# them; only the port check (run.py\textquotesingle{}s job) is skipped: Gunicorn binds the port.}
\CommentTok{\#}
\CommentTok{\# gunicorn arguments:}
\CommentTok{\#   "app:create\_app()"    — import \textquotesingle{}app\textquotesingle{}, call \textquotesingle{}create\_app()\textquotesingle{} to get the Flask app}
\CommentTok{\#   {-}{-}bind 0.0.0.0:5000   — listen on all interfaces (required in containers)}
\CommentTok{\#   {-}{-}workers 2           — 2 worker processes (handles 2 concurrent requests)}
\CommentTok{\#   {-}{-}timeout 30          — kill workers that take longer than 30 seconds}
\CommentTok{\#   {-}{-}access{-}logfile {-}    — write access logs to stdout (captured by Docker)}
\CommentTok{\#   {-}{-}error{-}logfile {-}     — write Gunicorn\textquotesingle{}s own messages to stderr (captured by Docker)}
\CommentTok{\#   {-}{-}log{-}level info      — log level for Gunicorn\textquotesingle{}s own messages}
\KeywordTok{CMD}\NormalTok{ [}\StringTok{"gunicorn"}\NormalTok{, }\OperatorTok{\textbackslash{}}
     \StringTok{"app:create\_app()"}\NormalTok{, }\OperatorTok{\textbackslash{}}
     \StringTok{"{-}{-}bind"}\NormalTok{, }\StringTok{"0.0.0.0:5000"}\NormalTok{, }\OperatorTok{\textbackslash{}}
     \StringTok{"{-}{-}workers"}\NormalTok{, }\StringTok{"2"}\NormalTok{, }\OperatorTok{\textbackslash{}}
     \StringTok{"{-}{-}timeout"}\NormalTok{, }\StringTok{"30"}\NormalTok{, }\OperatorTok{\textbackslash{}}
     \StringTok{"{-}{-}access{-}logfile"}\NormalTok{, }\StringTok{"{-}"}\NormalTok{, }\OperatorTok{\textbackslash{}}
     \StringTok{"{-}{-}error{-}logfile"}\NormalTok{, }\StringTok{"{-}"}\NormalTok{, }\OperatorTok{\textbackslash{}}
     \StringTok{"{-}{-}log{-}level"}\NormalTok{, }\StringTok{"info"}\NormalTok{]}
\end{Highlighting}
\end{Shaded}

\end{codeboxcap}

\noindent{}Add \texttt{gunicorn} to \texttt{requirements.txt} by hand, pinned like everything else:

\begin{outputbox}[short]
\begin{Verbatim}[fontsize=\codesize, breaklines, breakanywhere, breaksymbolleft={\tiny\textcolor{muted}{$\hookrightarrow$}}]
gunicorn==22.0.0
\end{Verbatim}
\end{outputbox}

\noindent{}Do not use \texttt{pip\ freeze\ \textgreater{}\ requirements.}\allowbreak{}\texttt{txt} here: your virtual environment now also contains the development tools from \hyperref[ch:17-build-process]{Chapter 17}, and freezing it would write pytest, black, ruff, and the rest into the production requirements~--- and into the image.

\subsection*{\texorpdfstring{The \texttt{.dockerignore} file}{The .dockerignore file}}\phantomsection\label{18-containerization:the-dockerignore-file}

\texttt{.dockerignore} tells Docker which files to exclude from the build context (the files sent to the Docker daemon). Without it, every file in the project~--- including \texttt{.git/}, \texttt{venv/}, \texttt{.env}, test files, and documentation~--- is included in the build context, slowing down builds and potentially including sensitive files in the image.

\begin{outputbox}
\begin{Verbatim}[fontsize=\codesize, breaklines, breakanywhere, breaksymbolleft={\tiny\textcolor{muted}{$\hookrightarrow$}}]
# .dockerignore
# Files and directories to exclude from the Docker build context.
# Syntax is the same as .gitignore.

# Python artifacts
venv/
__pycache__/
*.pyc
*.pyo
.pytest_cache/
.mypy_cache/
.ruff_cache/
htmlcov/
.coverage

# Environment files — NEVER include in the image
.env
.env.*
!.env.example

# Version control
.git/
.gitignore

# Development files
requirements-dev.txt
tests/
Makefile

# Documentation
*.md

# Editor and OS files
.DS_Store
*.swp
*.swo
.idea/
.vscode/
\end{Verbatim}
\end{outputbox}

\setcounter{section}{3}
\section{How Docker Layers Work}\label{18-containerization:184-how-docker-layers-work}

Understanding Docker layers is important for writing efficient Dockerfiles. Every Dockerfile instruction that produces filesystem changes creates a new \textbf{layer}. Layers are:

\begin{itemize}
\tightlist
\item
  \textbf{Cached}: Docker caches each layer. If a layer's instruction has not changed and all preceding layers are also unchanged, Docker reuses the cached layer~--- it does not re-run the instruction.
\item
  \textbf{Immutable}: once created, layers are read-only. They are shared between images and containers.
\item
  \textbf{Stacked}: the filesystem a container sees is all layers stacked together (using overlayfs).
\end{itemize}

\subsection*{Layer caching in practice}\phantomsection\label{18-containerization:layer-caching-in-practice}

A simplified, single-stage version of Listing 18.1 shows the principle:

\begin{codebox}[short]

\begin{Shaded}
\begin{Highlighting}[]
\CommentTok{\# Layer 1 — base image (cached after first pull)}
\KeywordTok{FROM}\NormalTok{ python:3.11{-}slim}

\CommentTok{\# Layer 2 — create the user (cached until the RUN command changes)}
\KeywordTok{RUN} \ExtensionTok{useradd} \AttributeTok{{-}{-}create{-}home} \AttributeTok{{-}{-}uid}\NormalTok{ 1001 notes}

\CommentTok{\# Layer 3 — requirements.txt copy (cached until requirements.txt changes)}
\KeywordTok{COPY}\NormalTok{ requirements.txt .}

\CommentTok{\# Layer 4 — pip install (cached until requirements.txt changes)}
\KeywordTok{RUN} \ExtensionTok{pip}\NormalTok{ install }\AttributeTok{{-}r}\NormalTok{ requirements.txt}

\CommentTok{\# Layer 5 — application code (cached until any app/ file changes)}
\KeywordTok{COPY}\NormalTok{ app/ ./app/}
\KeywordTok{COPY}\NormalTok{ run.py .}
\end{Highlighting}
\end{Shaded}

\end{codebox}

\noindent{}When you change only \texttt{app/}\allowbreak{}\texttt{routes/}\allowbreak{}\texttt{notes.py}:

\begin{itemize}
\tightlist
\item
  Layers 1--4: \textbf{cache hit} (reused, no re-execution)
\item
  Layer 5: \textbf{cache miss} (COPY app/ sees changed files, re-executes)
\end{itemize}

\noindent{}Build time: fraction of a second (only layer 5 is re-built).

If you had placed \texttt{COPY\ .\ .} before \texttt{pip\ install\ -}\allowbreak{}\texttt{r\ requirements.}\allowbreak{}\texttt{txt}:

\begin{itemize}
\tightlist
\item
  Any code change would invalidate the COPY layer
\item
  The pip install layer follows the COPY layer and is also invalidated
\item
  Build time: minutes (pip reinstalls all packages on every code change)
\end{itemize}

\noindent{}\textbf{The order of Dockerfile instructions matters enormously for build speed.} Always put things that change infrequently (base image, system packages, dependency files) before things that change frequently (application code).

\setcounter{section}{4}
\section{Building and Running the Image}\label{18-containerization:185-building-and-running-the-image}

\subsection*{Building the image}\phantomsection\label{18-containerization:building-the-image}

\begin{codebox}[short]

\begin{Shaded}
\begin{Highlighting}[]
\CommentTok{\# From the notes{-}api/ root directory}
\ExtensionTok{docker}\NormalTok{ build }\AttributeTok{{-}t}\NormalTok{ notes{-}api:latest .}

\CommentTok{\# {-}t notes{-}api:latest : tag the image with the name "notes{-}api" and version "latest"}
\CommentTok{\# .                   : use the current directory as the build context}
\end{Highlighting}
\end{Shaded}

\end{codebox}

\noindent{}First build output (no cache):

\begin{outputbox}[short]
\begin{Verbatim}[fontsize=\codesize, breaklines, breakanywhere, breaksymbolleft={\tiny\textcolor{muted}{$\hookrightarrow$}}]
[+] Building 21.4s (15/15) FINISHED
 => [internal] load build definition from Dockerfile
 => [internal] load .dockerignore
 => [builder 1/4] FROM python:3.11-slim
 => [builder 2/4] WORKDIR /build
 => [builder 3/4] COPY requirements.txt .
 => [builder 4/4] RUN pip install --no-cache-dir --prefix=/install -r requirements.txt
 => [final 1/7] FROM python:3.11-slim
 => [final 2/7] RUN useradd --create-home --shell /bin/bash --uid 1001 notes
 => [final 3/7] COPY --from=builder /install /usr/local
 => [final 4/7] WORKDIR /app
 => [final 5/7] RUN chown notes:notes /app
 => [final 6/7] COPY --chown=notes:notes app/ ./app/
 => [final 7/7] COPY --chown=notes:notes run.py .
 => exporting to image
\end{Verbatim}
\end{outputbox}

\noindent{}Second build (only code changed):

\begin{outputbox}[short]
\begin{Verbatim}[fontsize=\codesize, breaklines, breakanywhere, breaksymbolleft={\tiny\textcolor{muted}{$\hookrightarrow$}}]
[+] Building 1.3s (15/15) FINISHED
 => CACHED [builder 4/4] RUN pip install --no-cache-dir
 => CACHED [final 3/7] COPY --from=builder /install /usr/local
 => CACHED [final 5/7] RUN chown notes:notes /app
 => [final 6/7] COPY --chown=notes:notes app/ ./app/   ← only this layer re-ran
 => [final 7/7] COPY --chown=notes:notes run.py .      ← and this one
\end{Verbatim}
\end{outputbox}

\noindent{}1.3 seconds instead of 21. Layer caching is working.

\subsection*{Running the container with SQLite}\phantomsection\label{18-containerization:running-the-container-with-sqlite}

\begin{codebox}[short]

\begin{Shaded}
\begin{Highlighting}[]
\ExtensionTok{docker}\NormalTok{ run }\DataTypeTok{\textbackslash{}}
  \AttributeTok{{-}{-}name}\NormalTok{ notes{-}api }\DataTypeTok{\textbackslash{}}
  \AttributeTok{{-}{-}rm} \DataTypeTok{\textbackslash{}}
  \AttributeTok{{-}p}\NormalTok{ 5000:5000 }\DataTypeTok{\textbackslash{}}
  \AttributeTok{{-}e}\NormalTok{ DATABASE\_URL=}\StringTok{"sqlite:///notes.db"} \DataTypeTok{\textbackslash{}}
  \AttributeTok{{-}e}\NormalTok{ SECRET\_KEY=}\StringTok{"your{-}secret{-}key{-}here{-}at{-}least{-}24{-}chars"} \DataTypeTok{\textbackslash{}}
  \AttributeTok{{-}e}\NormalTok{ LOG\_FORMAT=human }\DataTypeTok{\textbackslash{}}
\NormalTok{  notes{-}api:latest}
\end{Highlighting}
\end{Shaded}

\end{codebox}

\noindent{}(\texttt{-e\ LOG\_FORMAT=}\allowbreak{}\texttt{human} overrides the image's JSON default, so the output below is easy to read.)

\textbf{Breaking down the \texttt{docker\ run} command:}

\Needspace{18\baselineskip}

{\def\LTcaptype{none} % do not increment counter
\begin{longtable}[]{@{}
  >{\raggedright\arraybackslash}p{(\linewidth - 2\tabcolsep) * \real{0.3067}}
  >{\raggedright\arraybackslash}p{(\linewidth - 2\tabcolsep) * \real{0.6933}}@{}}
\toprule\noalign{}
\begin{minipage}[b]{\linewidth}\raggedright
Flag
\end{minipage} & \begin{minipage}[b]{\linewidth}\raggedright
Meaning
\end{minipage} \\
\midrule\noalign{}
\endhead
\bottomrule\noalign{}
\endlastfoot
\texttt{-\/-name\ notes-api} & Name the container ``notes-api'' (instead of a random name) \\
\texttt{-\/-rm} & Remove the container when it stops (cleanup) \\
\texttt{-p\ 5000:5000} & Map host port 5000 to container port 5000 (\texttt{host:container}) \\
\texttt{-e\ DATABASE\_URL=}\allowbreak{}\texttt{...} & Set environment variable inside the container \\
\texttt{-e\ SECRET\_KEY=...} & Set another environment variable \\
\texttt{-e\ LOG\_FORMAT=}\allowbreak{}\texttt{human} & Override a default set by \texttt{ENV} in the Dockerfile \\
\texttt{notes-api:latest} & The image to run \\
\end{longtable}
}

After the container starts:

\begin{outputbox}[short]
\begin{Verbatim}[fontsize=\codesize, breaklines, breakanywhere, breaksymbolleft={\tiny\textcolor{muted}{$\hookrightarrow$}}]
[2024-01-15 14:00:00 +0000] [1] [INFO] Starting gunicorn 22.0.0
[2024-01-15 14:00:00 +0000] [1] [INFO] Listening at: http://0.0.0.0:5000 (1)
[2024-01-15 14:00:00 +0000] [1] [INFO] Using worker: sync
[2024-01-15 14:00:00 +0000] [7] [INFO] Booting worker with pid: 7
[2024-01-15 14:00:00 +0000] [8] [INFO] Booting worker with pid: 8
2024-01-15 14:00:00.412 INFO     app.logging_config | Logging configured. level=INFO format=human log_dir=(console only)
2024-01-15 14:00:00.413 INFO     app.startup        | Running startup checks...
2024-01-15 14:00:00.413 INFO     app.startup        | All startup checks passed.
2024-01-15 14:00:00.414 INFO     app                | Creating Flask application...
2024-01-15 14:00:00.431 INFO     app.database       | Database initialized.
2024-01-15 14:00:00.432 INFO     app                | Flask application created successfully.
... (the same six lines again, from the second worker)
\end{Verbatim}
\end{outputbox}

\noindent{}Gunicorn's master process (PID 1 in the container) starts first and binds the port. Each worker then imports the application and calls \texttt{create\_app()} itself~--- so the startup checks and \texttt{init\_db()} run once per worker, and every application line appears twice. (\texttt{CREATE\ TABLE\ IF\ NOT\ EXISTS} makes the repeat harmless. Gunicorn's \texttt{-\/-preload} option would create the app once, in the master, instead.)

Test it:

\begin{codebox}[short]

\begin{Shaded}
\begin{Highlighting}[]
\ExtensionTok{curl}\NormalTok{ http://localhost:5000/health}
\CommentTok{\# \{"status":"ok"\}}
\end{Highlighting}
\end{Shaded}

\end{codebox}

\subsection*{The port mapping explained}\phantomsection\label{18-containerization:the-port-mapping-explained}

\begin{diagrambox}[short]{197.3pt}
\begin{Verbatim}[fontsize=\fontsize{9.0}{8.55}\selectfont]
Your laptop             Container
┌──────────────┐        ┌──────────────┐
│              │        │              │
│  localhost   │        │  0.0.0.0     │
│  port 5000   │◄──────►│  port 5000   │
│              │  -p    │              │
│  (host)      │ 5000:  │  (container) │
│              │ 5000   │              │
└──────────────┘        └──────────────┘
\end{Verbatim}
\end{diagrambox}

\noindent{}The \texttt{-p\ 5000:5000} flag tells Docker to forward traffic arriving at the host's port 5000 to the container's port 5000. Without this flag, the container's port is not accessible from outside the container.

Inside the container, Gunicorn binds to \texttt{0.0.0.0:5000} (all interfaces). The \texttt{-\/-bind\ 0.0.0.0:5000} argument in the Dockerfile's \texttt{CMD} is what causes this~--- not \texttt{127.0.0.1} (loopback only, which inside a container means \emph{the container's own} loopback), but \texttt{0.0.0.0} (all interfaces, so Docker's port forwarding can reach it).

\setcounter{section}{5}
\section{Docker Compose: Running Multiple Containers Together}\label{18-containerization:186-docker-compose-running-multiple-containers-together}

The Notes API depends on a database. In development, it uses SQLite (a local file). In production, it uses PostgreSQL. Running both the application and PostgreSQL in Docker containers~--- and making them communicate~--- requires Docker Compose.

\textbf{Docker Compose} is a tool for defining and running multi-container applications using a YAML file.

\subsection*{\texorpdfstring{\texttt{docker-compose.yml}}{docker-compose.yml}}\phantomsection\label{18-containerization:docker-composeyml}

\begin{codeboxcap}{\textcolor{accent}{Listing 18.2}\enspace docker-compose.yml}

\begin{Shaded}
\begin{Highlighting}[]
\CommentTok{\# Define the multi{-}container development environment for the Notes API.}
\CommentTok{\#}
\CommentTok{\# Usage:}
\CommentTok{\#   docker compose up        — start all services}
\CommentTok{\#   docker compose up {-}d     — start in background (detached mode)}
\CommentTok{\#   docker compose down      — stop and remove containers}
\CommentTok{\#   docker compose logs {-}f   — follow logs from all services}
\CommentTok{\#   docker compose exec app bash  — open a shell inside the app container}
\CommentTok{\#}
\CommentTok{\# (Older Compose files start with a "version:" line. Compose v2 ignores it}
\CommentTok{\# and prints a warning, so it is left out.)}

\FunctionTok{services}\KeywordTok{:}

\CommentTok{  \# ───────────────────────────────────────────────────────────────────────────}
\CommentTok{  \# The PostgreSQL database service}
\CommentTok{  \# ───────────────────────────────────────────────────────────────────────────}
\AttributeTok{  }\FunctionTok{db}\KeywordTok{:}
\CommentTok{    \# Use the official PostgreSQL 15 image from Docker Hub}
\AttributeTok{    }\FunctionTok{image}\KeywordTok{:}\AttributeTok{ postgres:15{-}alpine}
\CommentTok{    \# \textquotesingle{}alpine\textquotesingle{} is a minimal Linux distribution — the image is roughly half the size}

\CommentTok{    \# Named volume for persistent storage.}
\CommentTok{    \# Without this, the database contents are lost when the container is}
\CommentTok{    \# removed or recreated (docker compose down, then up).}
\CommentTok{    \# With this, data persists in a Docker{-}managed volume.}
\AttributeTok{    }\FunctionTok{volumes}\KeywordTok{:}
\AttributeTok{      }\KeywordTok{{-}}\AttributeTok{ postgres\_data:/var/lib/postgresql/data}

\CommentTok{    \# Environment variables for the PostgreSQL container.}
\CommentTok{    \# POSTGRES\_USER: the database user to create}
\CommentTok{    \# POSTGRES\_PASSWORD: the password for that user}
\CommentTok{    \# POSTGRES\_DB: the database to create}
\AttributeTok{    }\FunctionTok{environment}\KeywordTok{:}
\AttributeTok{      }\FunctionTok{POSTGRES\_USER}\KeywordTok{:}\AttributeTok{ notes\_user}
\AttributeTok{      }\FunctionTok{POSTGRES\_PASSWORD}\KeywordTok{:}\AttributeTok{ notes\_password}
\AttributeTok{      }\FunctionTok{POSTGRES\_DB}\KeywordTok{:}\AttributeTok{ notesdb}

\CommentTok{    \# Health check for PostgreSQL.}
\CommentTok{    \# \textquotesingle{}pg\_isready\textquotesingle{} is a PostgreSQL utility that checks if the server is}
\CommentTok{    \# accepting connections. The app service will not start until this passes.}
\AttributeTok{    }\FunctionTok{healthcheck}\KeywordTok{:}
\AttributeTok{      }\FunctionTok{test}\KeywordTok{:}\AttributeTok{ }\KeywordTok{[}\StringTok{"CMD{-}SHELL"}\KeywordTok{,}\AttributeTok{ }\StringTok{"pg\_isready {-}U notes\_user {-}d notesdb"}\KeywordTok{]}
\AttributeTok{      }\FunctionTok{interval}\KeywordTok{:}\AttributeTok{ 5s}
\AttributeTok{      }\FunctionTok{timeout}\KeywordTok{:}\AttributeTok{ 5s}
\AttributeTok{      }\FunctionTok{retries}\KeywordTok{:}\AttributeTok{ }\DecValTok{5}

\CommentTok{    \# Expose PostgreSQL\textquotesingle{}s port only to other containers in this compose file.}
\CommentTok{    \# Do NOT use ports: here (which would expose it to the host) unless you}
\CommentTok{    \# need to connect with a local psql client.}
\AttributeTok{    }\FunctionTok{expose}\KeywordTok{:}
\AttributeTok{      }\KeywordTok{{-}}\AttributeTok{ }\StringTok{"5432"}

\CommentTok{  \# ───────────────────────────────────────────────────────────────────────────}
\CommentTok{  \# The Notes API application service}
\CommentTok{  \# ───────────────────────────────────────────────────────────────────────────}
\AttributeTok{  }\FunctionTok{app}\KeywordTok{:}
\CommentTok{    \# Build the image from the Dockerfile in the current directory.}
\CommentTok{    \# \textquotesingle{}context: .\textquotesingle{} means use the current directory as the build context.}
\AttributeTok{    }\FunctionTok{build}\KeywordTok{:}
\AttributeTok{      }\FunctionTok{context}\KeywordTok{:}\AttributeTok{ .}
\AttributeTok{      }\FunctionTok{dockerfile}\KeywordTok{:}\AttributeTok{ Dockerfile}

\CommentTok{    \# Port mapping: host port 5000 → container port 5000}
\AttributeTok{    }\FunctionTok{ports}\KeywordTok{:}
\AttributeTok{      }\KeywordTok{{-}}\AttributeTok{ }\StringTok{"5000:5000"}

\CommentTok{    \# Environment variables for the application.}
\CommentTok{    \# These override the ENV defaults set in the Dockerfile.}
\AttributeTok{    }\FunctionTok{environment}\KeywordTok{:}
\CommentTok{      \# Construct the DATABASE\_URL using the \textquotesingle{}db\textquotesingle{} service name as the hostname.}
\CommentTok{      \# Docker Compose creates a shared network where services can reach each}
\CommentTok{      \# other by their service name. \textquotesingle{}db\textquotesingle{} resolves to the PostgreSQL container\textquotesingle{}s}
\CommentTok{      \# IP address inside the Docker network.}
\AttributeTok{      }\FunctionTok{DATABASE\_URL}\KeywordTok{:}\AttributeTok{ postgresql://notes\_user:notes\_password@db:5432/notesdb}

\CommentTok{      \# In a real project, SECRET\_KEY would come from a secrets manager,}
\CommentTok{      \# not be hardcoded here. For development, this is acceptable.}
\AttributeTok{      }\FunctionTok{SECRET\_KEY}\KeywordTok{:}\AttributeTok{ dev{-}secret{-}key{-}change{-}in{-}production{-}12345}

\AttributeTok{      }\FunctionTok{FLASK\_DEBUG}\KeywordTok{:}\AttributeTok{ }\StringTok{"false"}
\AttributeTok{      }\FunctionTok{LOG\_FORMAT}\KeywordTok{:}\AttributeTok{ human}
\AttributeTok{      }\FunctionTok{LOG\_LEVEL}\KeywordTok{:}\AttributeTok{ INFO}

\CommentTok{    \# Wait for the database to be healthy before starting.}
\CommentTok{    \# \textquotesingle{}depends\_on\textquotesingle{} with \textquotesingle{}condition: service\_healthy\textquotesingle{} ensures Docker waits}
\CommentTok{    \# until the db healthcheck passes before starting the app.}
\AttributeTok{    }\FunctionTok{depends\_on}\KeywordTok{:}
\AttributeTok{      }\FunctionTok{db}\KeywordTok{:}
\AttributeTok{        }\FunctionTok{condition}\KeywordTok{:}\AttributeTok{ service\_healthy}

\CommentTok{    \# Restart policy: if the container crashes, restart it.}
\CommentTok{    \# \textquotesingle{}on{-}failure\textquotesingle{} restarts only if the container exits with a non{-}zero code.}
\AttributeTok{    }\FunctionTok{restart}\KeywordTok{:}\AttributeTok{ on{-}failure}

\CommentTok{\# Named volumes: Docker manages these on disk.}
\CommentTok{\# They persist between \textquotesingle{}docker compose down\textquotesingle{} and \textquotesingle{}docker compose up\textquotesingle{}.}
\CommentTok{\# To delete them: docker compose down {-}v}
\FunctionTok{volumes}\KeywordTok{:}
\AttributeTok{  }\FunctionTok{postgres\_data}\KeywordTok{:}
\end{Highlighting}
\end{Shaded}

\end{codeboxcap}

\subsection*{Running with Docker Compose}\phantomsection\label{18-containerization:running-with-docker-compose}

\begin{codebox}[short]

\begin{Shaded}
\begin{Highlighting}[]
\CommentTok{\# Start both services}
\ExtensionTok{docker}\NormalTok{ compose up}

\CommentTok{\# Expected output (abridged):}
\CommentTok{\# [+] Running 2/2}
\CommentTok{\#  ✔ Container notes{-}api{-}db{-}1   Started}
\CommentTok{\#  ✔ Container notes{-}api{-}app{-}1  Started}
\CommentTok{\#}
\CommentTok{\# notes{-}api{-}db{-}1   | ... database system is ready to accept connections}
\CommentTok{\# notes{-}api{-}app{-}1  | [...] [1] [INFO] Starting gunicorn 22.0.0}
\CommentTok{\# notes{-}api{-}app{-}1  | [...] [7] [INFO] Booting worker with pid: 7}
\CommentTok{\# notes{-}api{-}app{-}1  | [...] [8] [INFO] Booting worker with pid: 8}
\CommentTok{\# notes{-}api{-}app{-}1  | ... INFO     app.startup        | Running startup checks...}
\CommentTok{\# notes{-}api{-}app{-}1  | ... INFO     app.startup        | Database connectivity: OK}
\CommentTok{\# notes{-}api{-}app{-}1  | ... INFO     app.startup        | All startup checks passed.}
\CommentTok{\# notes{-}api{-}app{-}1  | ... INFO     app.database       | Database initialized.}
\CommentTok{\# (each application line appears once per worker)}
\end{Highlighting}
\end{Shaded}

\end{codebox}

\noindent{}Test it:

\begin{codebox}[short]

\begin{Shaded}
\begin{Highlighting}[]
\CommentTok{\# Create a note}
\ExtensionTok{curl} \AttributeTok{{-}X}\NormalTok{ POST http://localhost:5000/notes/ }\DataTypeTok{\textbackslash{}}
  \AttributeTok{{-}H} \StringTok{"Content{-}Type: application/json"} \DataTypeTok{\textbackslash{}}
  \AttributeTok{{-}d} \StringTok{\textquotesingle{}\{"content": "Running in a container!"\}\textquotesingle{}}

\CommentTok{\# \{"content":"Running in a container!","created\_at":"2024{-}01{-}15T14:00:05.123456+00:00","id":1\}}

\CommentTok{\# Retrieve it}
\ExtensionTok{curl}\NormalTok{ http://localhost:5000/notes/}
\CommentTok{\# \{"count":1,"notes":[\{"content":"Running in a container!","created\_at":"2024{-}01{-}15T14:00:05.123456+00:00","id":1\}]\}}
\end{Highlighting}
\end{Shaded}

\end{codebox}

\noindent{}Stop:

\begin{codebox}[short]

\begin{Shaded}
\begin{Highlighting}[]
\ExtensionTok{docker}\NormalTok{ compose down}
\end{Highlighting}
\end{Shaded}

\end{codebox}

\noindent{}The data persists in the \texttt{postgres\_data} volume. Running \texttt{docker\ compose\ up} again brings it back. (Flask's \texttt{jsonify} sorts keys alphabetically, which is why \texttt{content} comes before \texttt{id}; PostgreSQL's timestamps carry microseconds, SQLite's do not.)

\setcounter{section}{6}
\section{Image Tagging and Versioning}\label{18-containerization:187-image-tagging-and-versioning}

The tag \texttt{latest} is convenient but dangerous for deployment: it changes on every build, making it impossible to know which exact version is running on a server.

For production, tag images with a specific version identifier:

\begin{codebox}[short]

\begin{Shaded}
\begin{Highlighting}[]
\CommentTok{\# Tag with git commit hash — uniquely identifies the exact code version}
\VariableTok{GIT\_HASH}\OperatorTok{=}\VariableTok{$(}\FunctionTok{git}\NormalTok{ rev{-}parse }\AttributeTok{{-}{-}short}\NormalTok{ HEAD}\VariableTok{)}
\ExtensionTok{docker}\NormalTok{ build }\AttributeTok{{-}t}\NormalTok{ notes{-}api:sha{-}}\VariableTok{$\{GIT\_HASH\}} \AttributeTok{{-}t}\NormalTok{ notes{-}api:latest .}

\CommentTok{\# Example: notes{-}api:sha{-}a1b2c3d}
\end{Highlighting}
\end{Shaded}

\end{codebox}

\noindent{}Using the git commit hash as the image tag creates an unambiguous link between the running container and the exact commit it was built from. When a production incident occurs, you know the exact code version running. The \texttt{sha-} prefix is the convention this book uses everywhere (it is also what the CI tooling in \hyperref[ch:19-artifact-storage]{Chapter 19} produces), and it keeps a tag like \texttt{sha-1234567} from being mistaken for a version number. Tags are always lowercase. Deploy only these commit tags~--- never \texttt{latest}.

\subsection*{Adding version info to the image}\phantomsection\label{18-containerization:adding-version-info-to-the-image}

Record build metadata in the image using labels:

\begin{codebox}[short]

\begin{Shaded}
\begin{Highlighting}[]
\CommentTok{\# Add to Dockerfile, near the top of the final stage}
\KeywordTok{ARG}\NormalTok{ GIT\_HASH=unknown}
\KeywordTok{ARG}\NormalTok{ BUILD\_DATE=unknown}

\KeywordTok{LABEL}\NormalTok{ org.opencontainers.image.revision=}\StringTok{"$\{GIT\_HASH\}"} \OperatorTok{\textbackslash{}}
\NormalTok{      org.opencontainers.image.created=}\StringTok{"$\{BUILD\_DATE\}"} \OperatorTok{\textbackslash{}}
\NormalTok{      org.opencontainers.image.title=}\StringTok{"notes{-}api"} \OperatorTok{\textbackslash{}}
\NormalTok{      org.opencontainers.image.description=}\StringTok{"Notes REST API"}
\end{Highlighting}
\end{Shaded}

\end{codebox}

\noindent{}Build with the metadata:

\begin{codebox}[short]

\begin{Shaded}
\begin{Highlighting}[]
\VariableTok{GIT\_HASH}\OperatorTok{=}\VariableTok{$(}\FunctionTok{git}\NormalTok{ rev{-}parse }\AttributeTok{{-}{-}short}\NormalTok{ HEAD}\VariableTok{)}
\VariableTok{BUILD\_DATE}\OperatorTok{=}\VariableTok{$(}\FunctionTok{date} \AttributeTok{{-}u}\NormalTok{ +}\StringTok{"\%Y{-}\%m{-}\%dT\%H:\%M:\%SZ"}\VariableTok{)}

\ExtensionTok{docker}\NormalTok{ build }\DataTypeTok{\textbackslash{}}
  \AttributeTok{{-}{-}build{-}arg}\NormalTok{ GIT\_HASH=}\VariableTok{$\{GIT\_HASH\}} \DataTypeTok{\textbackslash{}}
  \AttributeTok{{-}{-}build{-}arg}\NormalTok{ BUILD\_DATE=}\VariableTok{$\{BUILD\_DATE\}} \DataTypeTok{\textbackslash{}}
  \AttributeTok{{-}t}\NormalTok{ notes{-}api:sha{-}}\VariableTok{$\{GIT\_HASH\}} \DataTypeTok{\textbackslash{}}
\NormalTok{  .}
\end{Highlighting}
\end{Shaded}

\end{codebox}

\noindent{}Inspect the metadata later:

\begin{codebox}[short]

\begin{Shaded}
\begin{Highlighting}[]
\ExtensionTok{docker}\NormalTok{ inspect }\AttributeTok{{-}{-}format} \StringTok{\textquotesingle{}\{\{json .Config.Labels\}\}\textquotesingle{}}\NormalTok{ notes{-}api:sha{-}a1b2c3d}
\CommentTok{\# \{"org.opencontainers.image.created":"2024{-}01{-}15T14:00:00Z",}
\CommentTok{\#  "org.opencontainers.image.description":"Notes REST API",}
\CommentTok{\#  "org.opencontainers.image.revision":"a1b2c3d",}
\CommentTok{\#  "org.opencontainers.image.title":"notes{-}api"\}}
\end{Highlighting}
\end{Shaded}

\end{codebox}

\setcounter{section}{7}
\section{Integrating Docker into the Build Process}\label{18-containerization:188-integrating-docker-into-the-build-process}

Add Docker build targets to the \texttt{Makefile}:

\begin{codebox}

\begin{Shaded}
\begin{Highlighting}[]
\CommentTok{\# Additions to Makefile — Docker targets}

\DataTypeTok{GIT\_HASH} \CharTok{:=}\StringTok{ }\CharTok{$(}\KeywordTok{shell}\StringTok{ git rev{-}parse {-}{-}short HEAD 2\textgreater{}/dev/null || echo "unknown"}\CharTok{)}
\DataTypeTok{BUILD\_DATE} \CharTok{:=}\StringTok{ }\CharTok{$(}\KeywordTok{shell}\StringTok{ date {-}u +"\%Y{-}\%m{-}\%dT\%H:\%M:\%SZ"}\CharTok{)}
\DataTypeTok{IMAGE\_NAME} \CharTok{:=}\StringTok{ notes{-}api}
\DataTypeTok{IMAGE\_TAG} \CharTok{:=}\StringTok{ sha{-}}\CharTok{$(}\DataTypeTok{GIT\_HASH}\CharTok{)}

\OtherTok{.PHONY:}\DataTypeTok{ docker{-}build docker{-}run docker{-}compose{-}up docker{-}compose{-}down docker{-}clean release}

\CommentTok{\# Build the Docker image}
\DecValTok{docker{-}build:}
\ErrorTok{    }\CharTok{@}\FunctionTok{echo }\StringTok{"→ Building Docker image }\CharTok{$(}\DataTypeTok{IMAGE\_NAME}\CharTok{)}\StringTok{:}\CharTok{$(}\DataTypeTok{IMAGE\_TAG}\CharTok{)}\StringTok{..."}
\NormalTok{    docker build \textbackslash{}}
        \CharTok{{-}}\FunctionTok{{-}build{-}arg GIT\_HASH=}\CharTok{$(}\DataTypeTok{GIT\_HASH}\CharTok{)}\FunctionTok{ }\CharTok{\textbackslash{}}
\FunctionTok{        {-}{-}build{-}arg BUILD\_DATE=}\CharTok{$(}\DataTypeTok{BUILD\_DATE}\CharTok{)}\FunctionTok{ }\CharTok{\textbackslash{}}
\FunctionTok{        {-}t }\CharTok{$(}\DataTypeTok{IMAGE\_NAME}\CharTok{)}\FunctionTok{:}\CharTok{$(}\DataTypeTok{IMAGE\_TAG}\CharTok{)}\FunctionTok{ }\CharTok{\textbackslash{}}
\FunctionTok{        {-}t }\CharTok{$(}\DataTypeTok{IMAGE\_NAME}\CharTok{)}\FunctionTok{:latest }\CharTok{\textbackslash{}}
\FunctionTok{        .}
    \CharTok{@}\FunctionTok{echo }\StringTok{"✓ Image built: }\CharTok{$(}\DataTypeTok{IMAGE\_NAME}\CharTok{)}\StringTok{:}\CharTok{$(}\DataTypeTok{IMAGE\_TAG}\CharTok{)}\StringTok{"}

\CommentTok{\# Run the container with SQLite (quick local test)}
\DecValTok{docker{-}run:}
\ErrorTok{    }\NormalTok{docker run {-}{-}rm \textbackslash{}}
        \CharTok{{-}}\FunctionTok{p 5000:5000 }\CharTok{\textbackslash{}}
\FunctionTok{        {-}e DATABASE\_URL=}\StringTok{"sqlite:///notes.db"}\FunctionTok{ }\CharTok{\textbackslash{}}
\FunctionTok{        {-}e SECRET\_KEY=}\StringTok{"docker{-}test{-}secret{-}key{-}1234"}\FunctionTok{ }\CharTok{\textbackslash{}}
\FunctionTok{        }\CharTok{$(}\DataTypeTok{IMAGE\_NAME}\CharTok{)}\FunctionTok{:latest}

\CommentTok{\# Start the full stack (app + PostgreSQL)}
\DecValTok{docker{-}compose{-}up:}
\ErrorTok{    }\NormalTok{docker compose up {-}d}
    \CharTok{@}\FunctionTok{echo }\StringTok{"✓ Stack started. API available at http://localhost:5000"}

\CommentTok{\# Stop the full stack}
\DecValTok{docker{-}compose{-}down:}
\ErrorTok{    }\NormalTok{docker compose down}

\CommentTok{\# Complete release build: run all checks, then build the Docker image}
\DecValTok{release:}\DataTypeTok{ build docker{-}build}
\ErrorTok{    }\CharTok{@}\FunctionTok{echo }\StringTok{""}
    \CharTok{@}\FunctionTok{echo }\StringTok{"━━━━━━━━━━━━━━━━━━━━━━━━━━━━━━━━━━━━━━━━━━━━━━━━━"}
    \CharTok{@}\FunctionTok{echo }\StringTok{"✓ Release build complete."}
    \CharTok{@}\FunctionTok{echo }\StringTok{"  Image: }\CharTok{$(}\DataTypeTok{IMAGE\_NAME}\CharTok{)}\StringTok{:}\CharTok{$(}\DataTypeTok{IMAGE\_TAG}\CharTok{)}\StringTok{"}
    \CharTok{@}\FunctionTok{echo }\StringTok{"━━━━━━━━━━━━━━━━━━━━━━━━━━━━━━━━━━━━━━━━━━━━━━━━━"}

\CommentTok{\# Remove all Docker images for this project (every tag, not just latest)}
\DecValTok{docker{-}clean:}
\ErrorTok{    }\NormalTok{docker rmi {-}f }\CharTok{$$}\NormalTok{(docker images {-}q }\CharTok{$(}\DataTypeTok{IMAGE\_NAME}\CharTok{)}\NormalTok{) 2\textgreater{}/dev/null || true}
    \CharTok{@}\FunctionTok{echo }\StringTok{"✓ Docker images removed."}
\end{Highlighting}
\end{Shaded}

\end{codebox}

\noindent{}Now:

\begin{codebox}[short]

\begin{Shaded}
\begin{Highlighting}[]
\FunctionTok{make}\NormalTok{ release}
\end{Highlighting}
\end{Shaded}

\end{codebox}

\noindent{}Runs the complete pipeline: lint → typecheck → test → security → docker build. The output is a Docker image built from code that passed every check, tagged with the current git commit hash.

\setcounter{section}{8}
\section{Adding Docker Tests to CI}\label{18-containerization:189-adding-docker-tests-to-ci}

Update the GitHub Actions workflow to also test the Docker build:

\begin{codebox}

\begin{Shaded}
\begin{Highlighting}[]
\CommentTok{\# Additions to .github/workflows/build.yml}

\AttributeTok{  }\FunctionTok{docker}\KeywordTok{:}
\AttributeTok{    }\FunctionTok{name}\KeywordTok{:}\AttributeTok{ Docker Build Test}
\AttributeTok{    }\FunctionTok{runs{-}on}\KeywordTok{:}\AttributeTok{ ubuntu{-}latest}
\AttributeTok{    }\FunctionTok{needs}\KeywordTok{:}\AttributeTok{ build}\CommentTok{  \# Only run if the build job passes}

\AttributeTok{    }\FunctionTok{steps}\KeywordTok{:}
\AttributeTok{      }\KeywordTok{{-}}\AttributeTok{ }\FunctionTok{name}\KeywordTok{:}\AttributeTok{ Check out repository}
\AttributeTok{        }\FunctionTok{uses}\KeywordTok{:}\AttributeTok{ actions/checkout@v4}

\AttributeTok{      }\KeywordTok{{-}}\AttributeTok{ }\FunctionTok{name}\KeywordTok{:}\AttributeTok{ Build Docker image}
\FunctionTok{        run}\KeywordTok{: }\CharTok{|}
\NormalTok{          docker build \textbackslash{}}
\NormalTok{            {-}{-}build{-}arg GIT\_HASH=$(git rev{-}parse {-}{-}short HEAD) \textbackslash{}}
\NormalTok{            {-}{-}build{-}arg BUILD\_DATE=$(date {-}u +"\%Y{-}\%m{-}\%dT\%H:\%M:\%SZ") \textbackslash{}}
\NormalTok{            {-}t notes{-}api:ci{-}test \textbackslash{}}
\NormalTok{            .}

\AttributeTok{      }\KeywordTok{{-}}\AttributeTok{ }\FunctionTok{name}\KeywordTok{:}\AttributeTok{ Test Docker image starts correctly}
\FunctionTok{        run}\KeywordTok{: }\CharTok{|}
\NormalTok{          \# Start the container in the background}
\NormalTok{          docker run {-}d \textbackslash{}}
\NormalTok{            {-}{-}name notes{-}api{-}test \textbackslash{}}
\NormalTok{            {-}p 5000:5000 \textbackslash{}}
\NormalTok{            {-}e DATABASE\_URL="sqlite:///notes.db" \textbackslash{}}
\NormalTok{            {-}e SECRET\_KEY="ci{-}test{-}secret{-}key{-}not{-}for{-}production{-}1234" \textbackslash{}}
\NormalTok{            notes{-}api:ci{-}test}

\NormalTok{          \# Test the health endpoint, retrying while the container starts}
\NormalTok{          \# (a fixed "sleep 5" is flaky on a slow runner)}
\NormalTok{          curl {-}{-}fail {-}{-}retry 10 {-}{-}retry{-}delay 1 {-}{-}retry{-}connrefused \textbackslash{}}
\NormalTok{            http://localhost:5000/health}

\NormalTok{          \# Stop and remove the container}
\NormalTok{          docker stop notes{-}api{-}test}
\NormalTok{          docker rm notes{-}api{-}test}
\end{Highlighting}
\end{Shaded}

\end{codebox}

\setcounter{section}{9}
\section{The Complete Updated Project Structure}\label{18-containerization:1810-the-complete-updated-project-structure}

\begin{diagrambox}[short]{377.0pt}
\begin{Verbatim}[fontsize=\fontsize{9.0}{8.55}\selectfont]
notes-api/
├── .github/
│   └── workflows/
│       └── build.yml          ← CI (updated with Docker job)
├── app/
│   ├── __init__.py
│   ├── config.py
│   ├── database.py
│   ├── errors.py
│   ├── logging_config.py
│   ├── middleware.py
│   ├── startup.py
│   └── routes/
│       ├── __init__.py
│       ├── health.py
│       └── notes.py
├── tests/
│   ├── __init__.py
│   ├── conftest.py
│   ├── test_errors.py
│   ├── test_health.py
│   └── test_notes.py
├── Dockerfile                 ← Multi-stage Docker image definition (NEW)
├── .dockerignore              ← Files excluded from Docker build (NEW)
├── docker-compose.yml         ← Multi-container development setup (NEW)
├── Makefile                   ← Build automation (updated with Docker targets)
├── pyproject.toml
├── requirements.txt           ← Updated with gunicorn
├── requirements-dev.txt
├── run.py
├── .env
├── .env.example
├── .python-version
└── .gitignore
\end{Verbatim}
\end{diagrambox}

\phantomsection\label{18-containerization:key-takeaways}
\begin{takeaways}

\begin{itemize}
\tightlist
\item
  A \textbf{container} is a Linux process with three kernel mechanisms: namespaces (isolation), cgroups (resource limits), and a layered filesystem (shared, immutable layers). Containers share the host kernel~--- they are not virtual machines.
\item
  \textbf{Docker} is the toolchain for building and running containers. A \textbf{Dockerfile} describes how to build an image. A \textbf{Docker image} is the immutable artifact. A \textbf{container} is a running instance of an image.
\item
  \textbf{Multi-stage builds} keep the final image small and secure: a builder stage installs (and, if necessary, compiles) dependencies; the final stage copies only what is needed at runtime. Build tools such as a compiler never appear in the production image.
\item
  \textbf{Layer caching} makes builds fast. Copy \texttt{requirements.txt} and run \texttt{pip\ install} \emph{before} copying application code~--- then code changes do not trigger a full pip reinstall.
\item
  \textbf{Never run as root} inside a container. Create a dedicated non-root user, give it the application files with \texttt{COPY\ -\/-chown}, and switch to it with \texttt{USER}.
\item
  \textbf{\texttt{PYTHONUNBUFFERED=}\allowbreak{}\texttt{1}} ensures log output appears immediately, not when a buffer fills. Critical for seeing the last log line before a crash.
\item
  \textbf{Use Gunicorn} (not \texttt{app.run()}) as the container's CMD. Gunicorn is a production WSGI server that handles concurrent requests and manages workers.
\item
  \textbf{Docker Compose} defines multi-container setups (application + database) in a single YAML file. Services reach each other by service name (\texttt{db} resolves to the PostgreSQL container's IP).
\item
  \textbf{Tag images with the git commit hash} (e.g., \texttt{notes-api:sha-a1b2c3d}) so you always know exactly which code is running in any environment.
\item
  \textbf{\texttt{make\ release}} is the complete pipeline: tests pass → Docker image built → artifact ready to deploy.
\end{itemize}

\end{takeaways}

\unnumberedsection{false}{Exercises}\label{18-containerization:exercises}

\begin{enumerate}
\def\labelenumi{\arabic{enumi}.}
\tightlist
\item
  \textbf{Predict.} Change one line of \texttt{app/}\allowbreak{}\texttt{routes/}\allowbreak{}\texttt{notes.py} and rebuild the image. Which layers are rebuilt, and which come from the cache? Now change \texttt{requirements.txt} instead.
\item
  \textbf{Break and fix.} Remove \texttt{USER\ notes} from the Dockerfile, build, and run \texttt{docker\ exec\ notes-}\allowbreak{}\texttt{api\ id}. Why is this dangerous? Restore it.
\item
  \textbf{Extend.} Measure the image size with \texttt{docker\ images}, then try \texttt{python:3.11-alpine} as the base. Did it get smaller? Did everything still work?
\item
  \textbf{Explain.} What exactly is shared between two containers started from the same image, and what is not?
\end{enumerate}

\noindent{}Solutions and hints are in the companion repository, in \texttt{exercises/}\allowbreak{}\texttt{chapter-18.md}. Try each exercise before reading its solution: the point is the thinking, not the answer.

\unnumberedsection{false}{What's Next}\label{18-containerization:whats-next}

The Notes API is now packaged in a Docker image that runs the same way everywhere. But the image exists only on the machine that built it. How does it get from your laptop~--- or from CI~--- to the server that will run it? Where are images kept, how are they named, and how do you know which image is which a month from now?

\noindent{}→ \hyperref[ch:19-artifact-storage]{Chapter 19}~--- Artifact Storage
